\documentclass[../main.tex]{subfiles}

\begin{document}
\chapter{Lebesgue Measurable Functions}

We devote this chapter to the study of measurable functions in order to lay the foundation for the study of the Lebesgue integral.

All the functions considered in this chapter take values (codomain) in the extended real numbers, that is, the set $\mathbb{R} \cup \{-\infty, +\infty\}$. This would make the discussion of limits and convergence of functions more convenient. The domain of the functions is assumed to be subsets of $\mathbb{R}$ as usual.

\section{Sums, Products, and Compositions}
Given two functions $h,g$ defined on a measurable set $E$, we often abbreviate $h < g$ to mean that $h(x) < g(x)$ for all $x \in E$, a sequence $\{f_n\}$ of functions is said to increase provided that $f_n < f_{n+1}$ for all $n \in \mathbb{N}$.

\begin{proposition}{}{}
	Let the function $f$ have a measurable domain $E$. Then the following statements are equivalent:
	\begin{itemize}
		\item For every $c \in \mathbb{R}$, the set $\{x \in E : f(x) > c\}$ is measurable.
		\item For every $c \in \mathbb{R}$, the set $\{x \in E : f(x) < c\}$ is measurable.
		\item For every $c \in \mathbb{R}$, the set $\{x \in E : f(x) \geq c\}$ is measurable.
		\item For every $c \in \mathbb{R}$, the set $\{x \in E : f(x) \leq c\}$ is measurable.
	\end{itemize}
	and each of these statements implies that the set $\{x \in E : f(x) = c\}$ is measurable for every $c \in \mathbb{R}$.
\end{proposition}
\begin{proof}
	We will show that the first statement implies the third, and the rest follows from taking complements. Just noting
	\begin{equation*}
		\{x \in E : f(x) \geq c\} = \bigcap_{n=1}^{\infty} \{x \in E : f(x) > c - \frac{1}{n}\}
	\end{equation*}
	would suffice.

	Next we show that the first statement implies $\{x \in E : f(x) = c\}$ is measurable. If $c\in \mathbb{R}$ we have
	\begin{equation*}
		\{x \in E : f(x) = c\} = \{x \in E : f(x) \geq c\} \cap \{x \in E : f(x) \leq c\}
	\end{equation*}
	so it is measurable. If $c = +\infty$ we have
	\begin{equation*}
		\{x \in E : f(x) = +\infty\} = \bigcap_{n=1}^{\infty} \{x \in E : f(x) > n\}
	\end{equation*}
	so it is measurable. If $c = -\infty$ we follow similar reasoning.
\end{proof}

\begin{definition}{Lebesgue Measurable Functions}{Lebesgue Measurable Functions}
	An extended real-valued function $f$ defined on a measurable set $E$ is said to be \textbf{Lebesgue measurable} if it satisfies any of the equivalent conditions in the previous proposition.
\end{definition}

\begin{example}{Characterization Function}{Characterization Function}
	Let $E$ be a measurable set and $A \subseteq E$. The characteristic function $\chi_A$ of $A$ is defined by
	\begin{equation*}
		\chi_A(x) =
		\begin{cases}
			1 & x \in A    \\
			0 & x \notin A
		\end{cases}
	\end{equation*}
	Then $\chi_A$ is measurable if and only if $A$ is measurable. This is because $\{x \in E : \chi_A(x) > c\}$ is either $A$, $E$, or $\emptyset$, depending on the value of $c$.
\end{example}

As opens sets in the extended real numbers can be countably generated by intervals, this naturally gives the following equivalent definition:

\begin{proposition}{Preimage of Open Sets is Measurable}{Preimage of Open Sets is Measurable}
	A function $f$ defined on a measurable set $E$ is Lebesgue measurable if and only if the preimage of every open set in the extended real numbers is measurable.
\end{proposition}

\begin{remark}
	In the context of extended real-valued functions, the topology of $\overline{\mathbb{R}}$ is generated by the open intervals $(a, b)$, $[-\infty, b)$, and $(a, \infty]$. Thus, an open set in $\overline{\mathbb{R}}$ may contain the points $\pm \infty$. This ensures that the preimage of a neighborhood of infinity, such as $\{x : f(x) > c\} = f^{-1}((c, \infty])$, is treated consistently within the definition of measurability.

	The topology of the extended real numbers $\overline{\mathbb{R}}$ contains the standard topology of $\mathbb{R}$. A set $U \subseteq \mathbb{R}$ is open in the standard sense if and only if it is open when viewed as a subset of $\overline{\mathbb{R}}$. This ensures that any Lebesgue measurable function $f: E \to \overline{\mathbb{R}}$ also behaves correctly when restricted to the set of points where it takes finite values.
\end{remark}

The following proposition tells us that the most familiar functions from elementary analysis, the continuous functions, are measurable.

\begin{proposition}{Continuous Functions are Measurable}{Continuous Functions are Measurable}
	A real-valued function that is continuous on its measurable domain is measurable.
\end{proposition}
\begin{proof}
	This is because the preimage of an open set under a continuous function is open, and hence measurable.
\end{proof}

\begin{proposition}{Monotone Functions are Measurable}{Monotone Functions are Measurable}
	A real-valued function that is monotone on its measurable domain is measurable.
\end{proposition}
\begin{proof}
	This is because $f^{-1}((c, \infty))$ is an interval, and hence measurable.
\end{proof}

\begin{proposition}{}{}
	Let $f$ be an extended real-valued function defined on a measurable set $E$.
	\begin{itemize}
		\item If $f$ is measurable on $E$, and $f = g$ a.e.\ on $E$, then $g$ is measurable on $E$.
		\item For a measurable subset $D \subseteq E$, $f$ is measurable on $E$ if and only if $f$ is measurable on $D$ and on $E \setminus D$.
	\end{itemize}
\end{proposition}
\begin{proof}
	Naturally, we need to separate the cases and define $A = \{x \in E : f(x) \neq g(x)\}$, and then consider $A$ and $E \setminus A$ separately. We have
	\begin{equation*}
		\{x \in E : g(x) > c\} = \{x \in A : g(x) > c\} \cup \{x \in E \setminus A : g(x) > c\}
	\end{equation*}
	and
	\begin{equation*}
		\{x \in E\setminus A : g(x) > c\} = \{x \in E \setminus A : f(x) > c\} = \{x \in E : f(x) > c\} \cap (E \setminus A)
	\end{equation*}
	which is measurable. Since $A$ is a zero measure set, $\{x \in A : g(x) > c\}$ being a subset of a zero measure set is also measurable. Hence, $\{x \in E : g(x) > c\}$ is measurable, and $g$ is measurable on $E$.

	For the second part, we assume $f$ is measurable on $E$. Then we have
	\begin{equation*}
		\{x \in D : f(x) > c\} = \{x \in E : f(x) > c\} \cap D
	\end{equation*}
	which is measurable, so does $\{x \in E \setminus D : f(x) > c\}$. Conversely, if $f$ is measurable on $D$ and $E \setminus D$, then we have
	\begin{equation*}
		\{x \in E : f(x) > c\} = \{x \in D : f(x) > c\} \cup \{x \in E \setminus D : f(x) > c\}
	\end{equation*}
	which is measurable, and $f$ is measurable on $E$.
\end{proof}

The sum $f+g$ of two measurable extended real-valued functions $f$ and $g$ is not well-defined if $f$ and $g$ take the values $+\infty$ and $-\infty$ at the same point.

If $f$ and $g$ are finite a.e.\ on $E$, we define $E_0 \subseteq E$ to be the set of points where $f$ and $g$ are both finite. Then $E_0$ is measurable, as
\begin{equation*}
	E_0 = \{x \in E : f(x) \in \mathbb{R}\} \cap \{x \in E : g(x) \in \mathbb{R}\}
\end{equation*}

If the restriction of $f+g$ to $E_0$ is measurable, as $E \setminus E_0$ is a zero measure set, we can apply any extended real-valued function to $E \setminus E_0$ and the resulting function will be measurable on $E$. This is the sense in which we consider it unambiguous to state that the sum of two measurable functions that are finite a.e.\ is measurable.

Similar remarks apply to products. The following proposition tells us that standard algebraic operations performed on measurable functions that are finite a.e.\ again lead to measurable functions.

\begin{theorem}{Algebra of Measurable Functions}{Algebra of Measurable Functions}
	Let $f$ and $g$ be measurable extended real-valued functions defined on a measurable set $E$, and they are finite a.e.\ on $E$. Then we have
	\begin{itemize}
		\item \textbf{Linearity:} $\forall \alpha, \beta \in \mathbb{R}$ the function $\alpha f + \beta g$ is measurable on $E$.
		\item \textbf{Product:} The function $fg$ is measurable on $E$.
	\end{itemize}
\end{theorem}
\begin{proof}
	It remains to show that $\alpha f$ is measurable for any $\alpha \in \mathbb{R}$. If $\alpha = 0$, then $\alpha f$ is the zero function, which is measurable. If $\alpha > 0$, then a scaling would do. So it suffice to consider the case $\alpha = \beta = 1$. We write
	\begin{equation}
		\{x \in E : f(x) + g(x) < c\} = \bigcup_{q \in \mathbb{Q}} \{x \in E : f(x) < q\} \cap \{x \in E : g(x) < c - q\}
	\end{equation}
	which is a countable union of measurable sets, and hence measurable.

	For the product, we write
	\begin{equation}
		fg = \frac{1}{2} \left[(f + g)^2 - f^2 - g^2\right]
	\end{equation}
	so it suffices to show that $f^2$ is measurable. But we have
	\begin{equation}
		\{x \in E : f^2(x) < c\} = \{x \in E : -\sqrt{c} < f(x) < \sqrt{c}\}
	\end{equation}
	which is measurable, and we are done.
\end{proof}

Many of the properties of functions considered in elementary analysis, including continuity and differentiability, are preserved under the operation of composition of functions. However, the composition of measurable functions may not be measurable.

\begin{example}{Non-measurable Compositions}{Non-measurable Compositions}
	There are two measurable real-valued functions, each defined on all of $\mathbb{R}$, whose composition fails to be measurable.

	In the previous chapter, we constructed a continuous, strictly increasing function $\psi$ on $[0,1]$ and a measurable subset $A \subseteq [0,1]$ such that $\psi(A)$ is not measurable. Extend $\psi$ to a continuous, strictly increasing function on all of $\mathbb{R}$, and $\psi^{-1}$ is also continuous and strictly increasing on all of $\mathbb{R}$, thus is measurable.

	On the other hand, the characteristic function $\chi_A$ of $A$ is measurable. Then the composition $\chi_A \circ \psi^{-1}$ is not measurable, since $\chi_A \circ \psi^{-1}(x) \geq 1$ if and only if $x \in \psi(A)$, which is not measurable.
\end{example}

Despite the setback imposed by this example, there is the following useful proposition regarding the preservation of measurability under composition.

\begin{proposition}{}{}
	Let $g$ be a measurable real-valued function defined on a measurable set $E$, and let $f$ be a continuous real-valued function defined on all of $\mathbb{R}$. Then the composition $f \circ g$ is measurable on $E$.
\end{proposition}
\begin{proof}
	Quite obvious as the preimage of an open set of a continuous function is open, and the preimage of an open set of a measurable function is measurable.
\end{proof}

\begin{remark}
	\textbf{Why Composition Fails?}

	The root of the failure in the composition of measurable functions lies in the difference between \textbf{Borel} and \textbf{Lebesgue} measurability.
	\begin{itemize}
		\item If $f$ were \textbf{Borel measurable}, then the preimage of \textit{any} Borel set would be Borel. Since every continuous function $g$ pulls back open sets to open sets (and thus Borel to Borel), the composition of two Borel functions is always Borel.
		\item However, for \textbf{Lebesgue measurable} functions, we only guarantee that the preimage of an \textbf{open} set is measurable. We do \textbf{not} guarantee that the preimage of a Lebesgue measurable set is measurable.
	\end{itemize}
	In the example $\chi_A \circ \psi^{-1}$, the set $A$ is Lebesgue measurable, but its preimage under the continuous function $\psi^{-1}$ is $\psi(A)$, which is non-measurable. The continuous function ``distorted'' the measure-zero set $A$ into a non-measurable ``monster'' $\psi(A)$.
\end{remark}

An immediate important consequence of the above composition result is that if $f$ is a measurable function on a measurable set $E$, then $|f|$ is also measurable, and every $|f|^p$ is measurable for $p > 0$.

Finite maxima and minima of measurable functions are also measurable, we denote
\begin{equation*}
	\max\{f,g\}(x) = \max\{f(x), g(x)\}, \quad \min\{f,g\}(x) = \min\{f(x), g(x)\}
\end{equation*}

\begin{proposition}{Finite Extrema Functions}{Finite Extrema Functions}
	For a finite family of measurable functions $\{f_1, f_2, \ldots, f_n\}$ defined on a measurable set $E$, the functions $\max\{f_1, f_2, \ldots, f_n\}$ and $\min\{f_1, f_2, \ldots, f_n\}$ are measurable on $E$.
\end{proposition}
\begin{proof}
	We have
	\begin{equation*}
		\{x \in E : \max\{f_1, f_2, \ldots, f_n\}(x) > c\} = \bigcup_{i=1}^{n} \{x \in E : f_i(x) > c\}
	\end{equation*}
	which is a finite union of measurable sets, and hence measurable. The proof for the minimum is similar.
\end{proof}

We also denote
\begin{equation*}
	|f|(x) = \max\{f(x), -f(x)\}, \quad f^+(x) = \max\{f(x), 0\}, \quad f^-(x) = \max\{-f(x), 0\}
\end{equation*}
for convenience, being the absolute value, positive part, and negative part of $f$, respectively. Then we have $f = f^+ - f^-$ and $|f| = f^+ + f^-$.

\section{Sequential Pointwise Limits and Simple Approximations}
For a sequence of functions $\{f_n\}$ defined on some set $E$ and a function $f$ defined on $E$, there are several ways to define what it means for $\{f_n\}$ to converge to $f$ on $E$. In this chapter we consider the concepts of pointwise convergence and uniform convergence, which are familiar from elementary analysis. In later chapters we consider many other modes of convergence for a sequence of functions.

\begin{definition}{Convergence}{Convergence}
	For a sequence of functions $\{f_n\}$ defined on some set $E$ and a function $f$ defined on $E$, and $A \subseteq E$, we say that
	\begin{itemize}
		\item $\{f_n\}$ converges to $f$ \textbf{pointwise} on $A$ if for every $x \in A$, $\lim_{n \to \infty} f_n(x) = f(x)$.
		\item $\{f_n\}$ converges to $f$ \textbf{a.e.\ pointwise} on $A$ if there exists a zero measure set $N \subseteq A$ such that $\{f_n\}$ converges to $f$ pointwise on $A \setminus N$.
		\item $\{f_n\}$ converges to $f$ \textbf{uniformly} on $A$ if for every $\epsilon > 0$, there exists an $N \in \mathbb{N}$ such that for all $n \geq N$ and all $x \in A$, $|f_n(x) - f(x)| < \epsilon$.
	\end{itemize}
\end{definition}

The pointwise limit of continuous functions may not be continuous. The pointwise limit of Riemann integrable functions may not be Riemann integrable. The following proposition is the first indication that the measurable functions have much better stability properties.

\begin{proposition}{}{}
	Let $\{f_n\}$ be a sequence of measurable functions on $E$ that converges pointwise a.e.\ on $E$ to a function $f$, then $f$ is measurable.
\end{proposition}
\begin{proof}
	We can just assume that $\{f_n\}$ converges pointwise on $E$ to $f$, otherwise replace $E$ with the subset of $E$ where the convergence occurs, and the remaining set is a zero measure set.

	To proceed, we note that for some $x$, $f(x) < c$ if and only if there exists a rational number $q < c$ such that $f(x) \leq q$, if and only if there exists a rational number $q < c$ and an $N \in \mathbb{N}$ such that $f_n(x) \leq q$ for all $n \geq N$. Thus we have
	\begin{equation*}
		\{x \in E : f(x) < c\} = \bigcup_{q \in \mathbb{Q}, q < c} \bigcup_{N=1}^{\infty} \bigcap_{n=N}^{\infty} \{x \in E : f_n(x) \leq q\}
	\end{equation*}
	which is a countable union of measurable sets, and hence measurable. Therefore, $f$ is measurable.
\end{proof}

Just like the step function being the building block of the Riemann integral, the simple function is the building block of the Lebesgue integral. A simple function is a linear combination of characteristic functions of measurable sets.

\begin{definition}{Simple Functions}{Simple Functions}
	A real-valued function $\varphi$ defined on a measurable set $E$ is called a \textbf{simple function} if it is measurable and takes only a finite number of values.
\end{definition}
in other words, a simple function can be expressed as
\begin{equation}
	\varphi(x) = \sum_{k=1}^{n} c_k \chi_{E_k}(x), \qquad E_k = \{x \in E : \varphi(x) = c_k\}
\end{equation}
being a finite linear combination of characteristic functions of measurable sets $E_k$. This expression is called \textbf{the canonical representation} of $\varphi$.

Simple functions are important because they are easy to understand and manipulate, and they can be used to approximate more complex measurable functions. In Riemann integration, we partition the domain into subintervals and approximate the function by a step function that takes constant values on each subinterval. In Lebesgue integration, we partition the range of the function into intervals and approximate the function by a simple function that takes constant values on the preimages of these intervals.

\begin{lemma}{The Simple Approximation Lemma}{The Simple Approximation Lemma}
	Let $f$ be a measurable real-valued function defined on a measurable set $E$. Assume $f$ is bounded on $E$. Then for every $\epsilon > 0$, there exists simple functions $\varphi_\epsilon$ and $\psi_\epsilon$ defined on $E$ such that
	\begin{equation}
		\varphi_\epsilon(x) \leq f(x) \leq \psi_\epsilon(x), \qquad 0 \leq \psi_\epsilon(x) - \varphi_\epsilon(x) < \epsilon, \quad \forall x \in E
	\end{equation}
\end{lemma}
\begin{proof}
	Let $(c,d)$ be an open interval containing $f(E)$, and we apply a partition
	\begin{equation*}
		c = y_0 < y_1 < \ldots < y_n = d
	\end{equation*}
	where $y_k - y_{k-1} < \epsilon$ for all $k$. Then we define $I_k = [y_{k-1}, y_k)$ for $k = 1, 2, \ldots, n-1$ and $E_k = f^{-1}(I_k)$, which is measurable as $f$ is measurable. Then we define
	\begin{equation}
		\varphi_\epsilon(x) = \sum_{k=1}^{n} y_{k-1} \chi_{E_k}(x), \qquad \psi_\epsilon(x) = \sum_{k=1}^{n} y_k \chi_{E_k}(x)
	\end{equation}
	which is the desired simple functions.
\end{proof}

\begin{theorem}{The Simple Approximation Theorem}{The Simple Approximation Theorem}
	An extended real-valued function $f$ on a measurable set $E$ is measurable if and only if there exists a sequence of simple functions $\{\varphi_n\}$ defined on $E$ such that it converges to $f$ pointwise on $E$ and
	\begin{equation}
		|\varphi_n(x)| \leq |f(x)|, \qquad \forall x \in E, n \in \mathbb{N}
	\end{equation}
	If $f$ is nonnegative, then we may choose $\{\varphi_n\}$ to be increasing.
\end{theorem}
\begin{proof}
	If $f$ is the limit of a sequence of simple functions $\{\varphi_n\}$, then $f$ is measurable as simple functions are measurable and the pointwise limit of measurable functions is measurable.

	Conversely, if $f$ is measurable, we first assume that $f \geq 0$. To handle the possibility that $f$ may be unbounded, we define $E_n = \{x \in E : f(x) < n\}$, which is measurable as $f$ is measurable. Then we apply the Simple Approximation Lemma to $f$ on $E_n$ with $\epsilon = 1/2^n$, and we obtain simple functions $\varphi_n$ and $\psi_n$ defined on $E_n$ such that
	\begin{equation}
		\varphi_n(x) \leq f(x) \leq \psi_n(x), \qquad 0 \leq \psi_n(x) - \varphi_n(x) < 1/2^n, \quad \forall x \in E_n
	\end{equation}
	Extend $\varphi_n$ to $E$ by defining $\varphi_n(x) = n$ for $x \in E \setminus E_n$. Then $\varphi_n$ is a simple function on $E$, and we have $0 \leq f(x) - \varphi_n(x) < 1/n$ for $x \in E_n$. It is very easy to see that $\varphi_n$ converges to $f$ pointwise on $E$: if $f(x)$ is finite, just take sufficiently large $n$ such that $x \in E_n$, and if $f(x) = +\infty$, then $\varphi_n(x) = n$ for all sufficiently large $n$, and hence $\varphi_n(x) \to +\infty = f(x)$.

	To make it increasing, we add a requirement that the partition of the interval $[0, n]$ is evenly spaced with interval length $1/2^n$, and hence is a refinement of the partition of the previous step. Then we have $\varphi_n \leq \varphi_{n+1}$ for all $n$.

	For the general case, we can write $f = f^+ - f^-$, and apply the above argument to $f^+$ and $f^-$ separately.
\end{proof}

\section{Littlewood's Three Principles, Egoroff's Theorem, and Lusin's Theorem}

Speaking of the theory of functions of a real variable, J. E. Littlewood says, ``The extent of knowledge required is nothing like so great as is sometimes supposed. There are three principles, roughly expressible in the following terms: Every [measurable] set is nearly a finite union of intervals; every [measurable] function is nearly continuous; every pointwise convergent sequence of [measurable] functions is nearly uniformly convergent. Most of the results of [the theory] are fairly intuitive applications of these ideas, and the student armed with them should be equal to most occasions when real variable theory is called for. If one of the principles would be the obvious means to settle the problem if it were `quite' true, it is natural to ask if the `nearly' is near enough, and for a problem that is actually solvable it generally is.''

We have already seen a formal formulation of Littlewood's first principle, \ref{thm:Littlewood First Principle}, that a measurable set with finite measure $E$, we know that for any $\epsilon > 0$, there is a finite disjoint collection of open intervals, whose union $\mathcal{U}$ is ``nearly'' equal to $E$ in the sense that $m(E \setminus \mathcal{U}) + m(\mathcal{U} \setminus E) < \epsilon$.

A precise realization of the last of Littlewood's principle is the following surprising theorem.

\begin{theorem}{Egoroff's Theorem}{Egoroff's Theorem}
	Assume that $E$ has finite measure, and let $\{f_n\}$ be a sequence of measurable functions on $E$ that converges pointwise on $E$ to a real-valued function $f$. Then for every $\epsilon > 0$, there exists a closed set $F \subseteq E$ such that $m(E \setminus F) < \epsilon$ and $\{f_n\}$ converges to $f$ uniformly on $F$.
\end{theorem}

\begin{remark}
	We have already done this frequently in basic analysis, where most questions require the usage of ``inner compact uniform convergence'' with dealing with problematic phenomenons at boundaries. For example, take $x^n$ on $[0,1]$, which converges pointwise to the function $f(x) = 0$ for $x \in [0,1)$ and $f(1) = 1$. This convergence is not uniform on $[0,1]$, but it is uniform on any closed interval $[0, 1 - \delta]$ for $\delta > 0$. This is a special case of Egoroff's theorem.
\end{remark}

To prove it, we first establish a lemma
\begin{proposition}{}{}
	Under the assumptions of Egoroff's Theorem, for each $\eta>0$ and $\delta > 0$, there exists a measurable subset $A \subseteq E$ and index $N$ such that $m(E \setminus A) < \delta$ and $|f_n(x) - f(x)| < \eta$ for all $x \in A$ and $n \geq N$.
\end{proposition}
\begin{remark}
	It may seem identical to the conclusion of Egoroff's theorem, but it is not. The difference is that the set $A$ may depend on $\eta$ and $\delta$, while in Egoroff's theorem, the set $F$ is independent of $\epsilon$.
\end{remark}
\begin{proof}
	Just take
	\begin{equation*}
		E_n = \{x \in E : |f_k(x) - f(x)| < \eta, \forall k \geq n\} = \bigcap_{k=n}^{\infty} \{x \in E : |f_k(x) - f(x)| < \eta\}
	\end{equation*}
	which is measurable as $f_n$ and $f$ are measurable. Then $E_n$ is increasing and $\bigcup_{n=1}^{\infty} E_n = E$, so we can take $N$ sufficiently large such that $m(E \setminus E_N) < \delta$, and then take $A = E_N$.
\end{proof}

\paragraph{Proof of Egoroff's Theorem.}
\begin{proof}
	The Lemma gives us a set where the sequence is ``nearly uniform'' for one fixed $\eta$. But for uniform convergence, you need the sequence to be ``nearly uniform'' for all $\eta$ at the same time. We can achieve this by applying the Lemma infinitely many times for smaller and smaller $\eta$, and then intersecting the results.

	Take $\eta_n = 1/n$ and $\delta_n = \epsilon/2^n$, and apply the Lemma to get a measurable set $A_n \subseteq E$ and index $N_n$ such that $m(E \setminus A_n) < \delta_n$ and $|f_k(x) - f(x)| < \eta_n$ for all $x \in A_n$ and $k \geq N_n$. Let $A = \bigcap_{n=1}^{\infty} A_n$. Then $A$ is measurable and
	\begin{equation*}
		m(E \setminus A) \leq \sum_{n=1}^{\infty} m(E \setminus A_n) < \sum_{n=1}^{\infty} \delta_n = \epsilon
	\end{equation*}
	Now, for any $\eta > 0$, choose sufficiently large $n$ such that $\eta_n < \eta$. Then for all $x \in A \subseteq A_n$ and $k \geq N_n$, we have $|f_k(x) - f(x)| < \eta_n < \eta$. This shows that $\{f_n\}$ converges to $f$ uniformly on $A$.

	If we further require $A$ to be closed, we can use the inner approximation property of measurable sets, just take a closed set $F \subseteq A$ such that $m(A \setminus F) < \epsilon$, then $m(E \setminus F) \leq m(E \setminus A) + m(A \setminus F) < 2\epsilon$. We can replace $\epsilon$ with $\epsilon/2$ in the above argument to get the desired result.
\end{proof}

\begin{remark}
	It is obvious that Egoroff's theorem also holds when the convergence is pointwise a.e.\ and the limit function is finite a.e. This is because we can just replace $E$ with the subset of $E$ where the convergence and finiteness occur, and the remaining set is a zero measure set.
\end{remark}

We now present a precise version of Littlewood's second principle in the case the measurable function is simple and then use this special case to prove the general case of the principle, Lusin's Theorem.

\begin{proposition}{}{}
	Let $f$ be a simple function defined on a measurable set $E$. Then for every $\epsilon > 0$, there exists a continuous function $g$ defined on $\mathbb{R}$ and a closed set $F \subseteq E$ such that $m(E \setminus F) < \epsilon$ and $f(x) = g(x)$ for all $x \in F$.
\end{proposition}
\begin{proof}
	We write the canonical representation of $f$ as
	\begin{equation*}
		f(x) = \sum_{k=1}^{n} a_k \chi_{E_k}(x), \qquad E_k = \{x \in E : f(x) = a_k\}
	\end{equation*}
	In each $E_k$, we can find a closed set $F_k \subseteq E_k$ such that $m(E_k \setminus F_k) < \epsilon/n$. Let $F = \bigcup_{k=1}^{n} F_k$, then $F$ is closed and $m(E \setminus F) < \epsilon$. Now we can define a continuous function $g$ on $\mathbb{R}$ such that $g(x) = f(x)$ for all $x \in F$.

	To extend $g$ to a continuous function on $\mathbb{R}$, consider $\mathbb{R} \setminus F$, which is open, can be expressed as a countable union of disjoint open intervals. On each of these intervals, we can define $g$ to be a linear interpolation between the values of $f$ on the endpoints of the interval (if they exist) or extend it to a constant value if the interval is unbounded. This ensures that $g$ is continuous on $\mathbb{R}$ and agrees with $f$ on the closed set $F$.
\end{proof}

\begin{remark}
	Actually, this need a more careful consideration here. We need to differ the case $f|_E$ is continuous and $f|_{\mathbb{R}}$ on $E$. The former take the subspace topology of $E$, and the latter take the standard topology of $\mathbb{R}$. As there may exists open sets in subspace topology that are not open in standard topology, we need to be careful.

	From the definitions, the subspace topology on $E$ is defined as taking intersection of open sets in $\mathbb{R}$ with $E$. So if an open set in $\mathbb{R}$ is already contained in $E$, then it is also open in the subspace topology. From this, we can say that
	\begin{itemize}
		\item If a function $f$ is continuous on $\mathbb{R}$, then $f|_E$ is continuous on $E$.
		\item If a function $f|_E$ is continuous on $E$, then $f$ may not be continuous on $\mathbb{R}$, take for example, the Dirichlet function and $E = \mathbb{Q}$, then $f|_E$ is continuous on $E$, but $f$ is nowhere continuous on $\mathbb{R}$.
	\end{itemize}

	Here, the limit of continuous uniformly convergent functions is continuous only gives the continuity of the limit function $g|_E$ is continuous on $E$, and actually we need to prove that $g$ is continuous on $E$. A general result here is called the \textbf{Tietze Extension Theorem}.

	A simplified proof is available for the case $X = \mathbb{R}$.
\end{remark}

\begin{proposition}{}{}
	Suppose $f$ is a continuous function defined on a closed set $F \subseteq \mathbb{R}$. Then there exists a continuous function $g$ defined on $\mathbb{R}$ such that $g(x) = f(x)$ for all $x \in F$.
\end{proposition}
\begin{proof}
	We have already constructed a function $g$ on $\mathbb{R}$ that agrees with $f$ on $F$ and is continuous on $\mathbb{R} \setminus F$. To show that $g$ is continuous on all of $\mathbb{R}$, we need to check the continuity at points in $F$.

	Take $x_0 \in F$, we know that $g|_F$ is continuous at $x_0$. We take any sequence of points $\{x_n\}$ in $\mathbb{R}$ that converges to $x_0$. And, take a subsequence that have one of the following cases, we try from the first case, and if one fails, we try the next one:
	\begin{itemize}
		\item If the subsequence is contained in $F$, then we're done, as $g|_F$ is continuous at $x_0$.
		\item If we cannot find a subsequence contained in $F$, then it seems that later in the sequence, all points are in $\mathbb{R} \setminus F$. Now, if we can find a subsequence contained in finitely many intervals of $\mathbb{R} \setminus F$, then $x_0$ must be an endpoint of one of these intervals, and a subsequence of this subsequence is contained in the interval, and we can use the interpolation to conclude that $g(x_n) \to g(x_0)$.
		\item If we cannot find a subsequence contained in finitely many intervals of $\mathbb{R} \setminus F$, then any neighborhood of $x_0$ contains infinitely many intervals of $\mathbb{R} \setminus F$. We find a subsequence that has one point in each of these intervals, and find a subsequence of this subsequence that is monotonically increasing or decreasing. Then, take two other sequences of points in $F$: one sequence is the bigger endpoint of the interval that contains the point in the subsequence, and the other sequence is the smaller endpoint of the interval that contains the point in the subsequence. Both of these sequences converge to $x_0$, and by the continuity of $g|_F$, we have $g$ evaluated at these sequences converges to $g(x_0)$. By the squeeze theorem, we conclude that $g(x_n) \to g(x_0)$.
	\end{itemize}
\end{proof}

\begin{theorem}{Lusin's Theorem}{Lusin's Theorem}
	Let $f$ be a real-valued measurable function defined on a measurable set $E$. Then for every $\epsilon > 0$, there exists a continuous function $g$ defined on $\mathbb{R}$ and a closed set $F \subseteq E$ such that $m(E \setminus F) < \epsilon$ and $f(x) = g(x)$ for all $x \in F$.
\end{theorem}
\begin{proof}
	First, assume that $m(E) < \infty$. By the Simple Approximation Theorem \ref{thm:Simple Approximation Theorem}, there exists a sequence of simple functions $\{\varphi_n\}$ defined on $E$ such that $\varphi_n \to f$ pointwise on $E$.

	Apply the previous proposition to each simple function $\varphi_n$ with $\epsilon/2^n$ to obtain a continuous function $g_n$ and a closed set $F_n \subseteq E$ such that $m(E \setminus F_n) < \epsilon/2^{n+1}$ and $\varphi_n(x) = g_n(x)$ for all $x \in F_n$.

	As we need to ensure that the limit of $\{g_n\}$ converges to a continuous function, we need uniform convergence. By the Egoroff's Theorem, we can find a closed set $F_0 \subseteq E$ such that $m(E \setminus F_0) < \epsilon/2$ and $\varphi_n \to f$ uniformly on $F_0$. This time, we can take $F = F_0 \cap \bigcap_{n=1}^{\infty} F_n$, which is closed and satisfies $m(E \setminus F) < \epsilon$.

	On $F$, we have $g_n = \varphi_n \to f$ uniformly, and since each $g_n$ is continuous, the uniform limit of continuous functions is continuous. Therefore, we can define $g = \lim_{n \to \infty} g_n$ on $F$, which is continuous on $F$ and agrees with $f$ on $F$. An extension of $g$ to all of $\mathbb{R}$ can be done similarly as in the previous proposition, ensuring that $g$ is continuous on $\mathbb{R}$. (The finite measure assumption is crucial here to ensure the validity of Egoroff's theorem.)

	\paragraph{The Infinite Measure Case ($m(E) = \infty$)}
	We partition the real line into a countable collection of disjoint, bounded intervals. For each $n \in \mathbb{Z}$, let $I_n = [n, n+1)$ and define $E_n = E \cap I_n$. Each $E_n$ is a measurable set of finite measure ($m(E_n) \leq 1$).

	Applying the result to each $E_n$, for a given $\epsilon > 0$, there exists a closed set $F_n \subseteq E_n$ such that:
	\begin{equation*}
    	m(E_n \setminus F_n) < \frac{\epsilon}{2^{|n|+2}} \quad \text{and} \quad f|_{F_n} \text{ is continuous.}
	\end{equation*}
	Define $F = \bigcup_{n \in \mathbb{Z}} F_n$. Since $F_n \subseteq E_n$, it follows that $F \subseteq E$. The measure of the excluded set is:
	\begin{equation*}
    	m(E \setminus F) = m\left( \bigcup_{n \in \mathbb{Z}} (E_n \setminus F_n) \right) = \sum_{n=-\infty}^{\infty} m(E_n \setminus F_n) < \sum_{n=-\infty}^{\infty} \frac{\epsilon}{2^{|n|+2}} = \epsilon.
	\end{equation*}

	We must show that $F$ is closed. Let $\{x_k\}$ be a sequence in $F$ such that $x_k \to x$. Since the sequence converges, it is bounded, meaning there exists some $M > 0$ such that $\{x_k\} \subseteq [-M, M]$. Consequently, the sequence $\{x_k\}$ can only intersect a finite number of the sets $\{F_n\}$. Because the union of a finite collection of closed sets is closed, and each $F_n$ is closed, the limit $x$ must belong to $\bigcup_{n=-M-1}^{M+1} F_n \subseteq F$. Thus, $F$ is closed.

	Furthermore, $f|_F$ is continuous. For any $x \in F$, $x$ belongs to some $F_n$. In a sufficiently small neighborhood of $x$, $f$ only interacts with finite many $F_k$. Since $f$ is continuous on each of these closed sets, $f|_F$ is continuous at $x$. 

	Finally, by the previously established extension proposition, there exists a continuous function $g$ defined on all of $\mathbb{R}$ such that $g(x) = f(x)$ for all $x \in F$. This completes the proof.
\end{proof}

\end{document}
